% ============================================================================
%  Documento de Diseño Técnico — Metaverso Escolar TecNM
%  Para el equipo de desarrollo. Compilar con: tectonic diseno-tecnico.tex
% ============================================================================
\documentclass[11pt,letterpaper]{article}

\usepackage[spanish,es-tabla]{babel}
\usepackage{geometry}
\geometry{letterpaper,top=2.3cm,bottom=2.3cm,left=2.4cm,right=2.4cm}
\usepackage{fontspec}
\usepackage{xcolor}
\usepackage{graphicx}
\usepackage{booktabs}
\usepackage{longtable}
\usepackage{array}
\usepackage{enumitem}
\usepackage{listings}
\usepackage{fancyhdr}
\usepackage{titlesec}
\usepackage{tikz}
\usetikzlibrary{shapes.geometric,arrows.meta,positioning,fit,backgrounds}
\usepackage[hidelinks]{hyperref}
\hypersetup{colorlinks=true,linkcolor=azultec,urlcolor=blue!60!black,
  pdftitle={Diseno Tecnico - Metaverso Escolar TecNM},pdfauthor={Daniel Octavio Ramirez Neri}}

\definecolor{azultec}{RGB}{31,59,110}
\definecolor{griscodigo}{RGB}{246,247,251}
\definecolor{verdecodigo}{RGB}{0,128,96}
\definecolor{rojo}{RGB}{176,0,32}

\setlength{\parskip}{0.5em}
\setlength{\parindent}{0pt}

\pagestyle{fancy}
\fancyhf{}
\fancyhead[L]{\small\itshape Diseño Técnico · Metaverso Escolar TecNM}
\fancyhead[R]{\small\thepage}
\renewcommand{\headrulewidth}{0.4pt}

\titleformat{\section}{\normalfont\Large\bfseries\color{azultec}}{\thesection}{0.6em}{}
\titleformat{\subsection}{\normalfont\large\bfseries}{\thesubsection}{0.5em}{}

\lstdefinestyle{codigo}{
  backgroundcolor=\color{griscodigo}, basicstyle=\ttfamily\small, breaklines=true,
  frame=single, rulecolor=\color{gray!40}, numbers=none,
  keywordstyle=\color{azultec}\bfseries, commentstyle=\itshape\color{gray},
  stringstyle=\color{verdecodigo}, showstringspaces=false, columns=fullflexible, tabsize=2, captionpos=b}
\lstset{style=codigo}

% Etiqueta de verbo HTTP
\newcommand{\http}[1]{\colorbox{azultec}{\textcolor{white}{\ttfamily\small\,#1\,}}}

\begin{document}

% ---- Portada simple ----
\begin{titlepage}
\centering
\vspace*{2.5cm}
{\Huge\bfseries\color{azultec} Documento de Diseño Técnico\par}
\vspace{0.6cm}
{\LARGE Metaverso Escolar TecNM\par}
\vspace{0.3cm}
{\large Backend, Panel Docente e Integración con Moodle (LTI 1.3)\par}
\vspace{2cm}
\begin{tabular}{rl}
\textbf{Versión:} & 1.0 \\[3pt]
\textbf{Fecha:} & Junio 2026 \\[3pt]
\textbf{Autor:} & Daniel Octavio Ramírez Neri \\[3pt]
\textbf{Audiencia:} & Equipo de desarrollo \\[3pt]
\textbf{Stack:} & Laravel 12 · PostgreSQL · Sanctum · LTI 1.3 · Unreal Engine \\
\end{tabular}
\vfill
{\small Este documento describe \emph{qué} se construyó, \emph{cómo} encaja y
\emph{por qué}, para que cualquier desarrollador pueda contribuir.\par}
\end{titlepage}

\tableofcontents
\clearpage

% ============================================================
\section{Visión general}
El \textbf{Metaverso Escolar} es un conjunto de \emph{prácticas} que el alumno
realiza dentro de un juego de Unreal Engine. El corazón del sistema es un
\textbf{backend} (Laravel + PostgreSQL) que administra la información académica y
expone una \textbf{API} que el juego consume. El alumno puede entrar de dos
formas: con un \textbf{magic link} (enlace de un solo uso) o desde una actividad
de \textbf{Moodle vía LTI 1.3} (que además devuelve la calificación al LMS).

\textbf{Tres consumidores, un cerebro:}
\begin{itemize}[leftmargin=1.3em,nosep]
  \item \textbf{Juego (Unreal)} $\rightarrow$ usa la API REST.
  \item \textbf{Panel docente (web)} $\rightarrow$ usa sesión web.
  \item \textbf{Moodle (LMS)} $\rightarrow$ se integra por LTI 1.3.
\end{itemize}

\subsection{Glosario rápido}
\begin{longtable}{p{3.2cm} p{12cm}}
\toprule \textbf{Término} & \textbf{Significado} \\ \midrule
\endhead
Magic link & Enlace con un token de un solo uso para lanzar el juego sin contraseña. \\
Sesión de práctica & Registro de una partida: estatus, calificación y telemetría. \\
LTI 1.3 & Estándar para conectar el LMS (Moodle) con una herramienta externa. \\
AGS & Servicio LTI para devolver calificaciones al LMS. \\
Deep Linking & Flujo LTI donde el docente elige \emph{qué} práctica enlazar. \\
JWKS & Conjunto de llaves públicas en JSON para validar firmas (JWT). \\
\bottomrule
\end{longtable}

% ============================================================
\section{Arquitectura}

\begin{figure}[h]
\centering
\begin{tikzpicture}[
  c/.style={rectangle,rounded corners,draw=azultec,thick,fill=azultec!8,
    text width=3.2cm,align=center,minimum height=1.1cm,font=\small},
  f/.style={-{Stealth[length=3mm]},thick,azultec}]
  \node[c] (moodle) {Moodle\\(plataforma LTI)};
  \node[c,right=2.2cm of moodle] (back) {Backend Laravel\\PostgreSQL · API};
  \node[c,below=1.3cm of back] (juego) {Juego Unreal};
  \node[c,below=1.3cm of moodle] (panel) {Panel docente};
  \draw[f] (moodle) -- node[above,font=\scriptsize]{LTI 1.3 / OIDC} (back);
  \draw[f] (back) -- node[right,font=\scriptsize]{API REST (Bearer)} (juego);
  \draw[f] (panel) -- node[left,font=\scriptsize]{sesión web} (back);
  \draw[f,dashed] (back) to[bend right=20] node[below,font=\scriptsize]{AGS: calificación} (moodle);
\end{tikzpicture}
\caption{Componentes y flujos principales.}
\end{figure}

\subsection{Stack tecnológico}
\begin{longtable}{p{4cm} p{11cm}}
\toprule \textbf{Componente} & \textbf{Tecnología / nota} \\ \midrule
\endhead
Backend & Laravel 12, PHP 8.3+ \\
Base de datos & PostgreSQL (usa \texttt{jsonb} para telemetría) \\
Auth de API & Laravel Sanctum (tokens Bearer con \emph{ability} \texttt{game}) \\
LTI & \texttt{packbackbooks/lti-1p3-tool} v6 \\
LMS & Moodle 5.0 (en Docker para desarrollo) \\
Cliente juego & Unreal Engine (C++), módulo HTTP nativo \\
Pruebas & Pest / PHPUnit (59 pruebas) \\
\bottomrule
\end{longtable}

\subsection{Estructura del repositorio (lo relevante)}
\begin{lstlisting}
app/
  Http/Controllers/Api/        # API del juego (redeem, sessions, complete, links, lti-redeem)
  Http/Controllers/Panel/      # Panel docente (login, dashboard, grupos, links, resultados)
  Http/Controllers/Lti/        # LTI (jwks, login, launch, deeplink)
  Http/Middleware/             # EnsurePanelAccess (middleware 'panel')
  Lti/                         # Adaptadores e implementaciones LTI (interfaces + Libreria*)
  Models/                      # Eloquent: Usuario, Alumno, Grupo, SesionPractica, LtiPlatform...
  Services/MagicLinkService.php
  Console/Commands/            # comandos artisan del proyecto
database/migrations/           # esquema completo (catalogo + operacion + LTI)
database/seeders/DemoSeeder.php
resources/views/{panel,lti}/   # vistas Blade
routes/{api.php,web.php}
tests/Feature/{Panel,Lti,...}/ # suite de pruebas
docs/                          # specs, planes, guias, este documento, ejemplo Unreal
moodle/docker-compose.yml      # Moodle local para pruebas LTI
\end{lstlisting}

% ============================================================
\section{Modelo de datos}
Esquema relacional normalizado, nombres en español, llaves primarias propias
(\texttt{id\_usuario}, \texttt{id\_grupo}, \ldots). Tablas clave:

\begin{longtable}{p{3.6cm} p{11cm}}
\toprule \textbf{Tabla} & \textbf{Notas importantes} \\ \midrule
\endhead
\texttt{usuarios} & \texttt{correo} único; \texttt{contrasena\_hash} oculto en serialización; \texttt{lti\_user\_id} para LTI. \\
\texttt{alumnos}/\texttt{maestros} & 1--1 con \texttt{usuarios}. \\
\texttt{grupos} & FK a materia, maestro y ciclo; clave \texttt{getRouteKeyName} = \texttt{id\_grupo}. \\
\texttt{practicas} & La actividad jugable; \texttt{escena\_referencia} (id de nivel en Unreal). \\
\texttt{eventos\_agenda} & Programa una práctica para un grupo. \\
\texttt{sesiones\_practica} & \texttt{estatus}, \texttt{calificacion} (float 0--100), \texttt{datos\_resultado} (\texttt{jsonb}); \texttt{id\_evento} \textbf{nullable} (las sesiones LTI se anclan a \texttt{id\_practica}); campos AGS. \\
\texttt{tokens\_juego} & \texttt{token\_hash} (SHA-256), \texttt{fecha\_expiracion}, \texttt{usado}, \texttt{ip\_origen}. \\
\texttt{lti\_platforms} & Moodle registrado (issuer, client\_id, deployment\_id, URLs). \\
\texttt{lti\_keys} & Par RSA de la herramienta (privada firma; pública en JWKS). \\
\texttt{lti\_nonces} & Soporte anti-replay (los nonces se manejan vía cache). \\
\bottomrule
\end{longtable}

\textbf{Decisiones de seguridad del esquema:} el token del magic link se guarda
\emph{hasheado} (nunca en claro); \texttt{id\_evento} se hizo nullable para no
forzar un evento de agenda en los lanzamientos por LTI.

% ============================================================
\section{Autenticación y seguridad}

\begin{longtable}{p{4.4cm} p{10.6cm}}
\toprule \textbf{Mecanismo} & \textbf{Detalle} \\ \midrule
\endhead
Magic link (juego) & Token aleatorio (64), se guarda solo su SHA-256; un solo uso (canje \textbf{atómico}); expiración configurable (\texttt{MAGIC\_LINK\_TTL\_MINUTES}, 120 por omisión). \\
Bearer del juego & Sanctum, \emph{ability} \texttt{game}; expira (\texttt{SANCTUM\_TOKEN\_TTL\_MINUTES}); requerido por \texttt{abilities:game} en las rutas de sesión. \\
Panel docente & Login por \textbf{URL firmada temporal} (\texttt{temporarySignedRoute}); middleware \texttt{panel} exige rol staff; rol Alumno bloqueado. \\
LTI & JWT firmado validado con el JWKS de Moodle; nonce de un solo uso; \texttt{/lti/*} excluido de CSRF (POST cross-site). \\
Datos sensibles & \texttt{contrasena\_hash} en \texttt{\$hidden}; nunca se serializa. \\
\bottomrule
\end{longtable}

% ============================================================
\section{Referencia de la API (juego)}
Prefijo \texttt{/api/game}. Respuestas JSON. Códigos de error consistentes:
\texttt{401} token inválido · \texttt{410} expirado/usado · \texttt{422} validación
· \texttt{403} recurso ajeno · \texttt{409} estado incorrecto.

\subsection{\http{POST} /api/game/redeem}
Canjea el magic link y entrega el Bearer.
\begin{lstlisting}
// request
{ "token": "<token-del-deeplink>" }
// 200
{ "access_token": "1|abc...", "token_type": "Bearer",
  "alumno": { "id_alumno": 1, "matricula": "20250001", ... },
  "practica": { "id_practica": 1, "titulo": "..." },
  "evento": { "id_evento": 1, "fecha_hora_inicio": "..." } }
\end{lstlisting}

\subsection{\http{POST} /api/game/lti-redeem}
Canje para sesiones creadas por LTI (devuelve la sesión ya existente).
\begin{lstlisting}
// request
{ "lti_session_token": "<token>" }
// 200
{ "access_token": "...", "token_type": "Bearer",
  "alumno": {...}, "practica": {...}, "id_sesion": 7 }
\end{lstlisting}

\subsection{Resto de endpoints (con Bearer + ability \texttt{game})}
\begin{longtable}{p{1.1cm} p{6cm} p{7cm}}
\toprule \textbf{Verbo} & \textbf{Ruta} & \textbf{Cuerpo / respuesta} \\ \midrule
\endhead
GET & \texttt{/api/game/me} & $\rightarrow$ \texttt{\{ usuario, alumno \}} \\
POST & \texttt{/api/game/sessions} & \texttt{\{ id\_evento \}} $\rightarrow$ \texttt{\{ id\_sesion, estatus \}} \\
POST & \texttt{/api/game/sessions/\{id\}/complete} & \texttt{\{ calificacion, datos\_resultado \}} $\rightarrow$ \texttt{200} \\
\bottomrule
\end{longtable}

Nota: \texttt{complete} valida propiedad (403 si la sesión es de otro alumno),
estado (409 si no está \texttt{en\_progreso}) y rango de \texttt{calificacion}
(0--100). Si la sesión es LTI, dispara el envío de la nota a Moodle (AGS).

% ============================================================
\section{Panel docente}
Páginas Blade bajo el middleware \texttt{panel}. Autorización por dueño del grupo
(Maestro ve lo suyo; Coordinador/Admin ven todo).

\begin{longtable}{p{1.1cm} p{6.6cm} p{6.4cm}}
\toprule \textbf{Verbo} & \textbf{Ruta} & \textbf{Función} \\ \midrule
\endhead
GET & \texttt{/panel/acceso/\{usuario\}} (firmada) & Inicia sesión del panel. \\
GET & \texttt{/panel} & Tablero de grupos. \\
GET & \texttt{/panel/grupos/\{grupo\}} & Alumnos y eventos del grupo. \\
POST & \texttt{/panel/grupos/\{g\}/eventos/\{e\}/links} & Genera magic links del grupo (+ CSV en cliente). \\
GET & \texttt{/panel/grupos/\{grupo\}/resultados} & Sesiones y calificaciones. \\
GET & \texttt{/panel/sesiones/\{sesion\}} & Detalle con telemetría. \\
\bottomrule
\end{longtable}

El inicio de sesión está encapsulado en un \textbf{único método}
(\texttt{PanelLoginController::establecerSesion}) para que LTI lo reutilice.

% ============================================================
\section{Integración LTI 1.3 (lo más importante para integrar)}

\subsection{Endpoints}
\begin{longtable}{p{1.1cm} p{4.6cm} p{8.4cm}}
\toprule \textbf{Verbo} & \textbf{Ruta} & \textbf{Función} \\ \midrule
\endhead
GET & \texttt{/lti/jwks} & Publica la llave pública de la herramienta. \\
GET/POST & \texttt{/lti/login} & Inicio OIDC; redirige a la plataforma. \\
POST & \texttt{/lti/launch} & Valida el lanzamiento; crea sesión o muestra el selector. \\
GET/POST & \texttt{/lti/deeplink} & Selector de práctica y respuesta firmada. \\
\bottomrule
\end{longtable}

\subsection{Desacople por interfaces (clave para pruebas)}
La lógica que toca la librería LTI está detrás de \textbf{interfaces} para que las
pruebas no dependan de criptografía real:
\begin{itemize}[leftmargin=1.3em,nosep]
  \item \texttt{LaunchValidador} $\rightarrow$ \texttt{LibreriaLaunchValidador}
  \item \texttt{DeepLinkRespondedor} $\rightarrow$ \texttt{LibreriaDeepLinkRespondedor}
  \item \texttt{AgsCliente} $\rightarrow$ \texttt{EnviarCalificacionAgs}
\end{itemize}
En pruebas se inyectan dobles; las implementaciones \texttt{Libreria*} se validan
en la prueba de extremo a extremo contra el Moodle real.

\subsection{Aprovisionamiento automático}
\texttt{AprovisionarAlumno::desdeLaunch()} busca el \texttt{Usuario} por
\texttt{lti\_user\_id}; si no existe lo crea (rol Alumno, correo determinístico si
falta) junto con su \texttt{Alumno}. Es \textbf{idempotente}.

\subsection{AGS — retorno de calificación}
Al completar una sesión LTI, \texttt{EnviarCalificacionAgs} mapea
\texttt{calificacion} a \texttt{scoreGiven} (con \texttt{scoreMaximum}=100) y la
envía al \emph{line item} de Moodle. El envío es resiliente: si Moodle no
responde, se registra el error sin romper el flujo del alumno.

\subsection{Configurar Moodle (resumen)}
1) Registrar la herramienta LTI 1.3 una vez (URLs de \texttt{/lti/login},
\texttt{/lti/launch}, \texttt{/lti/jwks}); 2) habilitar AGS y Deep Linking; 3)
correr \texttt{php artisan metaverso:lti-registrar-plataforma} con los datos que
Moodle genera. Detalle paso a paso en \texttt{docs/LTI-CONFIGURAR-MOODLE.md}.

\subsection{Notas de entorno local (¡leer antes de depurar!)}
Estos puntos resuelven los errores típicos de LTI sobre HTTP en Docker:
\begin{longtable}{p{5.2cm} p{9.4cm}}
\toprule \textbf{Síntoma} & \textbf{Causa / solución} \\ \midrule
\endhead
\textit{``cookies enabled''} en launch & Cookie \texttt{state} cross-site. Usar el navegador contra \texttt{localhost} (mismo dominio que Moodle). \\
\texttt{fix\_jwks\_alg(null)} & Moodle no leyó nuestra JWKS. Permitir el puerto del backend en \texttt{curlsecurityallowedport} y vaciar \texttt{curlsecurityblockedhosts} (bloquea IPs privadas como las de Docker). \\
\texttt{BeforeValidException} (iat) & Desfase de reloj del contenedor. Se aplica \texttt{JWT::\$leeway = 300} en \texttt{AppServiceProvider}. \\
``Invalid permissions'' (Moodle) & \texttt{moodledata} con dueño incorrecto. \texttt{chown -R daemon:daemon}. \\
\bottomrule
\end{longtable}

% ============================================================
\section{Cliente del juego (Unreal)}
El juego implementa este flujo mínimo:
\begin{enumerate}[leftmargin=1.5em,nosep]
  \item Recibe el token (del \emph{deep link} o por \texttt{-LtiToken=}).
  \item \texttt{POST /api/game/lti-redeem} $\rightarrow$ guarda \texttt{access\_token} y \texttt{id\_sesion}.
  \item Juega; al terminar: \texttt{POST /api/game/sessions/\{id\}/complete} con el Bearer y la calificación.
\end{enumerate}
Se entregó un \texttt{GameInstanceSubsystem} (C++) con
\texttt{IniciarSesion}, \texttt{SumarAcierto} y \texttt{Terminar}, invocables
desde Blueprint (ver \texttt{docs/unreal-ejemplo/}).

% ============================================================
\section{Cómo correr el proyecto localmente}
\begin{lstlisting}[language=bash]
# 1. Dependencias y entorno
composer install
cp .env.example .env && php artisan key:generate

# 2. Base de datos (PostgreSQL) + datos demo
createdb metaverso && createdb metaverso_test
php artisan migrate --seed

# 3. Servidor
php artisan serve --host=0.0.0.0 --port=8000

# 4. Pruebas
php artisan test            # 59 en verde

# 5. (Opcional) Moodle local para LTI
cd moodle && docker compose up -d   # http://localhost:8080 (admin / Admin12345!)
\end{lstlisting}

\subsection{Variables de entorno}
\begin{longtable}{p{5.6cm} p{3cm} p{6cm}}
\toprule \textbf{Variable} & \textbf{Default} & \textbf{Uso} \\ \midrule
\endhead
\texttt{MAGIC\_LINK\_TTL\_MINUTES} & 120 & Vigencia del magic link. \\
\texttt{GAME\_DEEPLINK\_SCHEME} & tecnm-metaverso & Esquema del deep link. \\
\texttt{SANCTUM\_TOKEN\_TTL\_MINUTES} & 480 & Vigencia del Bearer. \\
\texttt{LINKS\_API\_KEY} & (vacío) & Clave para \texttt{/api/links} (fail-closed). \\
\texttt{PANEL\_LOGIN\_TTL\_MINUTES} & 30 & Vigencia del enlace del panel. \\
\bottomrule
\end{longtable}

% ============================================================
\section{Pruebas}
La suite (Pest) cubre: flujo de magic link y canje (incluyendo 401/410 y
single-use), sesiones (start/complete con 403/409/422), panel (acceso firmado,
rol, autorización por grupo), y LTI (JWKS, aprovisionamiento, launch, deep
linking y AGS, con los dobles de interfaz). Total: \textbf{59 pruebas}. Ejecutar
con \texttt{php artisan test}. Las implementaciones reales de LTI se verifican en
la prueba de extremo a extremo contra Moodle.

% ============================================================
\section{Decisiones de diseño y trade-offs}
\begin{longtable}{p{5cm} p{9.6cm}}
\toprule \textbf{Decisión} & \textbf{Razón} \\ \midrule
\endhead
Token hasheado + un solo uso & Si se filtra la BD, no se regalan accesos; evita replay. \\
\texttt{id\_evento} nullable & Las sesiones LTI se anclan a una práctica, no a un evento de agenda. \\
Login del panel firmado & Sin contraseñas ni SMTP; un solo método reutilizable por LTI. \\
LTI tras interfaces & Pruebas deterministas sin criptografía; la librería se valida en e2e. \\
Auto-aprovisionar alumnos & Cero altas manuales: menos fricción para el docente. \\
AGS resiliente (no lanza) & Un fallo de Moodle no debe romper la partida del alumno. \\
\bottomrule
\end{longtable}

% ============================================================
\section{Trabajo futuro}
\begin{itemize}[leftmargin=1.3em,nosep]
  \item Producción: HTTPS, secretos gestionados, proteger \texttt{/api/links} con auth de docente.
  \item Panel: crear prácticas/eventos y mostrar sesiones LTI en resultados.
  \item LTI: NRPS (sincronización de listas) y multi-plataforma.
  \item Juego: contenido jugable y registro del esquema de deep link en el SO.
\end{itemize}

\vfill
\begin{center}\small\itshape
Documentos complementarios: \texttt{docs/COMO-FUNCIONA.md} (explicación sin
tecnicismos), \texttt{docs/INTEGRACION-MOODLE-LTI.md} (estrategia LTI),
\texttt{docs/LTI-CONFIGURAR-MOODLE.md} (configuración paso a paso),
\texttt{docs/superpowers/specs} y \texttt{plans} (specs y planes por sub-proyecto).
\end{center}

\end{document}
